\section{Dependent Local Tasks with Collaborative Actions}
\label{ch03:sec:dep-local-collab}

Robots must not only navigate the workspace, but also execute actions that
interact with the environment. This section first builds a unified model of
motion and actions, so that local temporal tasks can specify both where the
robot moves and which actions it performs. The resulting planning scheme extends
the nominal automata-based solution of Chapter~\ref{chap:preliminaries}.
The section then addresses dependent local tasks. Dependencies are represented
through collaborative and assisting actions rather than fixed robot identities.
Each robot requests assistance from neighbors, evaluates incoming requests,
confirms feasible collaborations, and adapts its local plan using only local
computation and communication. Since dependencies are formed and removed during
execution, the scheme remains scalable and resilient to robot failures.





%=============================================================================================

\subsection{Motion and action planning}
\label{ch03:sec:action-plan}

The specifications considered so far describe region visits. In many robotic
tasks, however, reaching a region enables operations whose feasibility depends
on regional properties, while executed actions modify internal robot states.
Motion and action therefore have to be modeled jointly. The construction below
separates workspace-dependent mobility information from workspace-independent
action knowledge, which improves modularity and reuse when the workspace or
specification changes \cite{van2008handbook}.

(I) \emph{Mobility abstraction}.
The robot mobility is modeled by the weighted finite transition system below:
\begin{equation}
\label{ch03:eq:M}
\mathcal{M} \triangleq
(\Pi_{\mathcal{M}}, \Sigma_{\mathcal{M}}, \longrightarrow_{\mathcal{M}},
 \Pi_{\mathcal{M},0}, \Psi_{\mathcal{M}}, L_{\mathcal{M}},
 W_{\mathcal{M}}),
\end{equation}
where $\Sigma_{\mathcal{M}} \triangleq \{\sigma_{\mathcal{M}}\}$ and the
structure follows Definition~\ref{ch02:def:FTS}. The set
$\Pi_{\mathcal{M}}$ contains workspace regions, $\Pi_{\mathcal{M},0}$ contains
initial regions, and
$\Psi_{\mathcal{M}} \triangleq \Psi_r \cup \Psi_p$ combines region names with
workspace properties relevant to actions. The action $\sigma_{\mathcal{M}}$ is
the symbolic counterpart of the continuous navigation controller in
\eqref{ch02:eq:control}.

(II) \emph{Action model}.
Action description languages encode preconditions and effects of task-level
actions \cite{gelfond1998action}, including STRIPS and ADL models
\cite{fikes1972strips,pednault1989adl}. Suppose the robot can perform
$K$ actions
$\{\sigma_{\mathcal{B},1}, \cdots, \sigma_{\mathcal{B},K}\}$ through
controllers $\{\mathcal{K}_1, \cdots, \mathcal{K}_K\}$, and define
$\Sigma_{\mathcal{B}} \triangleq
\{\sigma_{\mathcal{B},0}, \sigma_{\mathcal{B},1}, \cdots,
  \sigma_{\mathcal{B},K}\}$ with
$\sigma_{\mathcal{B},0} \triangleq \texttt{None}$. The set $\Psi_s$ represents
internal robot states, and $\Psi_b$ records the single active task-level action
under the non-concurrency assumption. Action enabledness is characterized by:
\begin{equation}
\label{ch03:eq:cond}
\texttt{Cond} :
\Sigma_{\mathcal{B}} \times 2^{\Psi_p} \times 2^{\Psi_s}
\longrightarrow \{\top,\bot\},
\end{equation}
where the value indicates whether an action is admissible under the current
workspace properties and internal state. By definition,
$\texttt{Cond}(\sigma_{\mathcal{B},0}, w_p, w_s)=\top$ for all
$w_p \subseteq \Psi_p$ and $w_s \subseteq \Psi_s$. The action-induced state
update is captured by:
\begin{equation}
\label{ch03:eq:eff}
\texttt{Eff} :
\Sigma_{\mathcal{B}} \times 2^{\Psi_s} \times \Psi_b
\longrightarrow 2^{\Psi_s} \times \Psi_b,
\end{equation}
where \eqref{ch03:eq:eff} updates the internal state valuation and the active
action indicator. The null action preserves the internal state and resets the
action status to \texttt{None}; other actions update the internal state and
activate the corresponding action proposition. Workspace properties are treated
as external inputs and are not modified by \texttt{Eff}.

\begin{definition}
Given $\Psi_p$, $\Psi_s$, $\Psi_b$, $\Sigma_{\mathcal{B}}$,
$\texttt{Cond}$ from \eqref{ch03:eq:cond}, and $\texttt{Eff}$ from
\eqref{ch03:eq:eff}, define the action map by:
\begin{equation}
\label{ch03:eq:B}
\mathcal{B} \triangleq
(\Pi_{\mathcal{B}}, \Sigma_{\mathcal{B}}, \Psi_p,
 \hookrightarrow_{\mathcal{B}}, \Pi_{\mathcal{B},0}, \Psi_{\mathcal{B}},
 L_{\mathcal{B}}, W_{\mathcal{B}}),
\end{equation}
where the individual components are specified below.
The state set is
$\Pi_{\mathcal{B}} \triangleq 2^{\Psi_s} \times \Psi_b$, and the initial state
set satisfies
$\Pi_{\mathcal{B},0} \subseteq 2^{\Psi_s} \times \{\Psi_{b,0}\}$. Moreover,
$\Psi_{\mathcal{B}} \triangleq \Psi_s \cup \Psi_b$. The propositions in
$\Psi_p$ play the role of external inputs, so the transition relation is
conditional. For
$\pi_{\mathcal{B}} = (w_s,\Psi_{b,\ell})$ and
$\pi_{\mathcal{B}}' = (w_s',\Psi_{b,k})$, the tuple
$(\pi_{\mathcal{B}}, \alpha_{\mathcal{B}}, w_p, \pi_{\mathcal{B}}')$ belongs
to $\hookrightarrow_{\mathcal{B}}$ if and only if
$\alpha_{\mathcal{B}} \in \Sigma_{\mathcal{B}}$,
$\texttt{Cond}(\alpha_{\mathcal{B}}, w_p, w_s)=\top$, and
$\texttt{Eff}(\alpha_{\mathcal{B}}, w_s, \Psi_{b,\ell})
= (w_s', \Psi_{b,k})$. The labeling function is given by
$L_{\mathcal{B}}(w_s,\Psi_{b,k}) \triangleq w_s \cup \{\Psi_{b,k}\}$, and
$W_{\mathcal{B}}$ assigns a positive weight to each enabled transition
according to the cost of the corresponding action.
\end{definition}


The state space of $\mathcal{B}$ is $2^{\Psi_s} \times \Psi_b$, rather than
$2^{\Psi_s} \times 2^{\Psi_b}$, because only one action proposition can be true.
The map is nondeterministic before workspace properties are fixed: any enabled
action may be selected. Once the external valuation is fixed, the conditional
transitions define a standard weighted finite transition system
\cite{baier2008principles}. Since $\mathcal{B}$ is independent of workspace
geometry, it can be reused across environments.

(III) \emph{Complete functionality model}.
The mobility abstraction in \eqref{ch03:eq:M} and the action map in
\eqref{ch03:eq:B} describe complementary aspects of robot behavior. Richer task
specifications involving both navigation and action execution require a unified
model.

\begin{definition}
The complete robot functionality is defined by the composition:
\begin{equation}
\label{ch03:eq:R}
\mathcal{R} \triangleq
(\Pi_{\mathcal{R}}, \Sigma_{\mathcal{R}}, \longrightarrow_{\mathcal{R}},
 \Pi_{\mathcal{R},0}, \Psi_{\mathcal{R}}, L_{\mathcal{R}},
 W_{\mathcal{R}}),
\end{equation}
with the components specified as follows.
The state set is
$\Pi_{\mathcal{R}} \triangleq \Pi_{\mathcal{M}} \times \Pi_{\mathcal{B}}$, and
the action set is
$\Sigma_{\mathcal{R}} \triangleq
\Sigma_{\mathcal{M}} \cup \Sigma_{\mathcal{B}}$.
The initial state set is
$\Pi_{\mathcal{R},0} \triangleq
\Pi_{\mathcal{M},0} \times \Pi_{\mathcal{B},0}$, and the proposition set is
$\Psi_{\mathcal{R}} \triangleq
\Psi_{\mathcal{M}} \cup \Psi_{\mathcal{B}}$.
The transition relation in \eqref{ch03:eq:R} has two types of transitions.
First, for a mobility action $\sigma_{\mathcal{M}}$, one has
$\langle \pi_{\mathcal{M}}, \pi_{\mathcal{B}} \rangle
\xrightarrow{\sigma_{\mathcal{M}}}_{\mathcal{R}}
\langle \pi_{\mathcal{M}}', \pi_{\mathcal{B}}' \rangle$
if
$\pi_{\mathcal{M}} \xrightarrow{\sigma_{\mathcal{M}}}_{\mathcal{M}}
\pi_{\mathcal{M}}'$
and
$\pi_{\mathcal{B}}' =
\texttt{Eff}(\sigma_{\mathcal{B},0}, w_s, \Psi_{b,\ell})$ for
$\pi_{\mathcal{B}} = (w_s,\Psi_{b,\ell})$. Second, for a task action
$\alpha_{\mathcal{B}} \in \Sigma_{\mathcal{B}}$, one has
$\langle \pi_{\mathcal{M}}, \pi_{\mathcal{B}} \rangle
\xrightarrow{\alpha_{\mathcal{B}}}_{\mathcal{R}}
\langle \pi_{\mathcal{M}}, \pi_{\mathcal{B}}' \rangle$
if
$(\pi_{\mathcal{B}}, \alpha_{\mathcal{B}},
  L_{\mathcal{M}}(\pi_{\mathcal{M}}) \cap \Psi_p,
  \pi_{\mathcal{B}}') \in \hookrightarrow_{\mathcal{B}}$.
Thus, a mobility transition changes the region, whereas a task transition keeps
the region fixed and updates the internal action state.
The labeling function is defined by
$L_{\mathcal{R}}(\langle \pi_{\mathcal{M}}, \pi_{\mathcal{B}} \rangle)
\triangleq
L_{\mathcal{M}}(\pi_{\mathcal{M}}) \cup L_{\mathcal{B}}(\pi_{\mathcal{B}})$.
The weight function combines the weights inherited from the two component
models. Mobility transitions use the weight assigned by $W_{\mathcal{M}}$,
while task transitions use the corresponding weight from $W_{\mathcal{B}}$.
\end{definition}


Note that during a mobility transition, the auxiliary action $\sigma_{\mathcal{B},0}$ is
applied automatically, meaning that no task-level action is active during
motion. The composition in \eqref{ch03:eq:R} attaches a copy of the action map
to each region, whose workspace properties determine the enabled action
transitions.

(IV) \emph{Local task over motion and action}.
The composed model $\mathcal{R}$ above is a weighted finite transition system over
$\Psi_{\mathcal{R}} = \Psi_r \cup \Psi_p \cup \Psi_s \cup \Psi_b$. Thus task
formulas can refer to locations, workspace properties, internal states, and
executed actions, enabling specifications over both action ordering and action
locations \cite{bhatia2010sampling,kloetzer2010automatic,smith2011optimal}. For
example, executing ``drop \textsf{A}'' at region \textsf{1} infinitely often is
expressed as $\varphi = \square \Diamond (\Psi_{r,1} \wedge \Psi_{b,2})$, while
preconditions enforce where \textsf{A} may be picked up.
Finally, the planning problem for hybrid control strategy
addressed in this subsection is stated as follows.

\begin{problem}
\label{ch03:pro:full}
Given the complete functionality model $\mathcal{R}$ and a task formula
$\varphi$, determine a motion-and-action plan that minimizes the cost in
\eqref{ch02:eq:prod-cost}, and construct a hybrid control strategy that
executes this plan.
\end{problem}


Since $\mathcal{R}$ in \eqref{ch03:eq:R} is a standard weighted finite
transition system, the synthesis procedure of Chapter~\ref{chap:preliminaries}
applies directly. Algorithm~\ref{ch02:alg:optimal} computes an optimal
accepting run $R_{\text{opt}}$ of
$\mathcal{R} \otimes \mathcal{A}_{\varphi}$, whose projection yields
$\tau_{\text{opt}} \triangleq R_{\text{opt}}|_{\Pi_{\mathcal{R}}}$. Each
transition of $\tau_{\text{opt}}$ is then implemented by the corresponding
low-level action or navigation controller, as in
Algorithm~\ref{ch02:alg:hybrid}. The framework therefore consists of building
$\mathcal{M}$ and $\mathcal{B}$, composing $\mathcal{R}$, synthesizing
$\tau_{\text{opt}}$, and executing it by hybrid control. Its modularity allows
$\mathcal{B}$ to be reused across workspaces, while changes in workspace or task
specification affect only the mobility abstraction or synthesis step.























%==============================
\subsection{Multi-robot systems with dependent local tasks}
\label{ch03:sec:collaborative}

The previous subsection considered independent local tasks. The next setting
allows the completion of one robot's task to require assistance from others, for
a heterogeneous team indexed by
$i \in \mathcal{N} \triangleq \{1,2,\cdots,N\}$.

(I) \emph{Mobility abstraction}.
The workspace is partitioned into $M$ regions of interest,
$\Pi \triangleq \{\pi_1,\pi_2,\cdots,\pi_M\}$. These symbols are assumed to be
fixed a priori and known to all robots. The motion of robot $i$ is modeled by
the weighted finite transition system:
\begin{equation}
\label{ch03:M}
\mathcal{M}_i \triangleq
(\Pi, \longrightarrow_{\mathcal{M}}^i, \Pi_{\mathcal{M},0}^i,
 \Psi_{\mathcal{M}}^i, L_{\mathcal{M}}^i, T_{\mathcal{M}}^i),
\end{equation}
where $\longrightarrow_{\mathcal{M}}^i \subseteq \Pi \times \Pi$ is the motion
transition relation, $\Pi_{\mathcal{M},0}^i \subseteq \Pi$ is the initial
region set, $L_{\mathcal{M}}^i : \Pi \rightarrow 2^{\Psi_{\mathcal{M}}^i}$ is
the labeling function, and
$T_{\mathcal{M}}^i : \longrightarrow_{\mathcal{M}}^i \rightarrow \mathbb{R}^+$
assigns a positive duration to each transition. A finite path of
$\mathcal{M}_i$ is a sequence $\pi_0 \pi_1 \cdots \pi_P$ such that
$(\pi_n,\pi_{n+1}) \in \longrightarrow_{\mathcal{M}}^i$ for all
$n = 0,1,\cdots,P-1$. Because of heterogeneity, the models $\mathcal{M}_i$ may
differ across the robots.
Each robot $i$ is also associated with a set of neighboring robots
$\mathcal{N}_i \subseteq \mathcal{N}$. Messages can be exchanged directly
between robot $i$ and any robot $j \in \mathcal{N}_i$.

(II) \emph{Action model}.
In addition to motion, robot $i$ has actions
$\Sigma_i \triangleq \Sigma_l^i \cup \Sigma_c^i \cup \Sigma_h^i$, consisting
of local actions, collaborative actions initiated by robot $i$, and assisting
actions that robot $i$ can provide. The active actions initiated by robot $i$
are $\Sigma_a^i \triangleq \Sigma_l^i \cup \Sigma_c^i$, with
$\sigma_0 \triangleq \texttt{None}$. Let $\Sigma_h^{\sim i}$ denote the
external assisting actions on which robot $i$ may depend, where
$\Sigma_h^{\sim i} \subseteq \bigcup_{j \in \mathcal{N}_i}\Sigma_h^j$. The
action model is
\begin{equation}
\label{ch03:A}
\mathcal{B}_i \triangleq
(\Sigma_i, \Psi_{\Sigma}^i, L_{\Sigma}^i, \texttt{Cond}_i,
 \texttt{Dura}_i, \texttt{Depd}_i),
\end{equation}
where $\Psi_{\Sigma}^i$ is the set of atomic propositions associated with the
active actions of robot $i$, and
$L_{\Sigma}^i : \Sigma_i \rightarrow 2^{\Psi_{\Sigma}^i}$ is the action
labeling function. For assisting actions,
$L_{\Sigma}^i(\sigma_h)=\emptyset$ for all $\sigma_h \in \Sigma_h^i$, while
$L_{\Sigma}^i(\sigma_a)\subseteq \Psi_{\Sigma}^i$ for all
$\sigma_a \in \Sigma_a^i$.
The function
$\texttt{Cond}_i : \Sigma_i \times 2^{\Psi_{\mathcal{M}}^i}
\rightarrow \{\top,\bot\}$ specifies the region properties required for an
action to be enabled. The function
$\texttt{Dura}_i : \Sigma_i \rightarrow \mathbb{R}^+$ assigns durations, with
$\texttt{Dura}_i(\sigma_0)=T_0>0$. The dependence function
$\texttt{Depd}_i : \Sigma_i \rightarrow 2^{\Sigma_h^{\sim i}}$ lists the
assisting actions required by each action. Local and assisting actions have
empty dependence sets, whereas collaborative actions may depend on elements of
$\Sigma_h^{\sim i}$. Because dependencies are defined over actions rather than
robot identities, the assisting robot need not be fixed in advance.

\begin{definition}
\label{ch03:done}
Let $\pi_v \in \Pi$. A local or assisting action
$\sigma_s \in \Sigma_l^i \cup \Sigma_h^i$ is said to be done at region $\pi_v$
if
$\texttt{Cond}_i(\sigma_s,L_{\mathcal{M}}^i(\pi_v))=\top$ and $\sigma_s$ is
activated for the duration $\texttt{Dura}_i(\sigma_s)$. A collaborative action
$\sigma_c \in \Sigma_c^i$ is said to be done at region $\pi_v$ if the same two
conditions hold and, in addition, every assisting action in
$\texttt{Depd}_i(\sigma_c)$ is done by other robots at the same region $\pi_v$.
\end{definition}

(III) \emph{Complete robot model}.
The complete model of robot $i$ combines its motion abstraction and its action
model into a single finite transition system.

\begin{definition}
\label{ch03:full-model}
Given $\mathcal{M}_i$ and $\mathcal{B}_i$, define the complete robot model by:
\begin{equation}
\label{ch03:eq:Ri}
\mathcal{R}_i \triangleq
(\Pi_{\mathcal{R}}^i, \longrightarrow_{\mathcal{R}}^i,
 \Pi_{\mathcal{R},0}^i, \Psi_{\mathcal{R}}^i,
 L_{\mathcal{R}}^i, T_{\mathcal{R}}^i),
\end{equation}
where the components are specified below.
The state set is
$\Pi_{\mathcal{R}}^i \triangleq \Pi \times \Sigma_i$, and each state is written
as $\pi_{\mathcal{R},s}=\langle \pi_t,\sigma_n \rangle$ with
$\pi_t \in \Pi$ and $\sigma_n \in \Sigma_i$. The initial state set is
$\Pi_{\mathcal{R},0}^i \triangleq
\{\langle \pi_{\mathcal{M},0}^i,\sigma_0 \rangle\}$, where
$\Pi_{\mathcal{M},0}^i=\{\pi_{\mathcal{M},0}^i\}$. The proposition set is
$\Psi_{\mathcal{R}}^i \triangleq \Psi_{\mathcal{M}}^i \cup \Psi_{\Sigma}^i$.
The transition relation
$\longrightarrow_{\mathcal{R}}^i \subseteq
\Pi_{\mathcal{R}}^i \times \Pi_{\mathcal{R}}^i$
is defined as follows. A pair
$\big(\langle \pi_s,\sigma_m \rangle,\langle \pi_t,\sigma_n \rangle\big)$
belongs to $\longrightarrow_{\mathcal{R}}^i$ if one of the following cases
holds. First, $\sigma_m=\sigma_n=\sigma_0$ and
$\pi_s \longrightarrow_{\mathcal{M}}^i \pi_t$, which corresponds to a motion
transition. Second, $\sigma_m=\sigma_0$, $\sigma_n \neq \sigma_0$,
$\pi_s=\pi_t$, and
$\texttt{Cond}_i(\sigma_n,L_{\mathcal{M}}^i(\pi_s))=\top$, which corresponds to
the start of an enabled action. Third,
$\sigma_n=\sigma_0$ and
$\sigma_m \in \Sigma_i \setminus \{\sigma_0\}$ with $\pi_s=\pi_t$, which
is the completion of the action.
The labeling function
$L_{\mathcal{R}}^i : \Pi_{\mathcal{R}}^i \rightarrow 2^{\Psi_{\mathcal{R}}^i}$
is defined by
$L_{\mathcal{R}}^i(\langle \pi_s,\sigma_m \rangle) \triangleq
L_{\mathcal{M}}^i(\pi_s) \cup L_{\Sigma}^i(\sigma_m)$. The transition duration
$T_{\mathcal{R}}^i : \longrightarrow_{\mathcal{R}}^i \rightarrow \mathbb{R}^+$
is given by the corresponding motion or action duration. Motion transitions
inherit the value of $T_{\mathcal{M}}^i$, action-start transitions use
$\texttt{Dura}_i(\sigma_n)$, and action-completion transitions use the fixed
duration $T_0$.
\end{definition}

The model $\mathcal{R}_i$ is a standard finite transition system
\cite{baier2008principles}. A finite path is written
$\tau_i=\pi_{\mathcal{R},0}\pi_{\mathcal{R},1}\cdots\pi_{\mathcal{R},P}$,
with trace
$\texttt{trace}(\tau_i)\triangleq
L_{\mathcal{R}}^i(\pi_{\mathcal{R},0})
L_{\mathcal{R}}^i(\pi_{\mathcal{R},1}) \cdots
L_{\mathcal{R}}^i(\pi_{\mathcal{R},P})$. The local task $\varphi_i$ is an
sc-safe LTL formula over $\Psi_{\mathcal{R}}^i$, and may constrain motion,
local actions, and collaborative actions. A path fulfills $\varphi_i$ if
$\texttt{trace}(\tau_i)\models \varphi_i$; the case
$\varphi_i \triangleq \top$ represents a pure assisting robot.

\begin{problem}
\label{ch03:big-problem}
Given $\mathcal{R}_i$ and the locally assigned task $\varphi_i$ of each robot
$i \in \mathcal{N}$, design a distributed control and coordination scheme such
that all tasks $\varphi_i$ are fulfilled, even when some local tasks require
collaboration from other robots.
\end{problem}

\subsection{Distributed task coordination}
\label{ch03:sec:coordination}

This subsection develops the bottom-up coordination scheme: local initial
synthesis, online request and reply, real-time plan adaptation, and recovery
from robot failures.

(I) \emph{Initial plan synthesis}.
Given $\mathcal{R}_i$ and $\varphi_i$, robot $i$ first synthesizes an initial
motion-and-action plan off-line. This plan is the starting point for the online
coordination mechanism. Let
$\mathcal{A}_{\varphi_i} \triangleq
(Q_i, 2^{\Psi_{\mathcal{R}}^i}, \delta_i, Q_{i,0}, \mathcal{F}_i)$
be the NBA associated with $\varphi_i$, as in
Section~\ref{ch02:sec:LTL}. The product automaton between $\mathcal{R}_i$ and
$\mathcal{A}_{\varphi_i}$ is given by:
\begin{equation}
\label{ch03:Ap}
\mathcal{A}_{p,i} \triangleq
\mathcal{R}_i \otimes \mathcal{A}_{\varphi_i}
= (Q_{p,i}, \delta_{p,i}, Q_{p,i,0}, \mathcal{F}_{p,i}, W_{p,i}),
\end{equation}
where
$Q_{p,i} \triangleq \Pi_{\mathcal{R}}^i \times Q_i$,
$Q_{p,i,0} \triangleq \Pi_{\mathcal{R},0}^i \times Q_{i,0}$, and
$\mathcal{F}_{p,i} \triangleq \Pi_{\mathcal{R}}^i \times \mathcal{F}_i$.
Moreover,
$\big(\langle \pi_{\mathcal{R},s},q_m \rangle,
\langle \pi_{\mathcal{R},k},q_n \rangle\big) \in \delta_{p,i}$
if
$\pi_{\mathcal{R},s} \longrightarrow_{\mathcal{R}}^i \pi_{\mathcal{R},k}$ and
$(q_m,L_{\mathcal{R}}^i(\pi_{\mathcal{R},s}),q_n) \in \delta_i$. The weight
function $W_{p,i} : \delta_{p,i} \rightarrow \mathbb{R}^+$ is defined by:
\begin{equation}
\label{ch03:eq:wp}
W_{p,i}\big(
\langle \pi_{\mathcal{R},s},q_m \rangle,
\langle \pi_{\mathcal{R},k},q_n \rangle
\big)
\triangleq
T_{\mathcal{R}}^i(\pi_{\mathcal{R},s},\pi_{\mathcal{R},k}),
\end{equation}
where the right-hand side uses the transition duration in $\mathcal{R}_i$.
There exists a finite path of $\mathcal{R}_i$ satisfying $\varphi_i$ if and
only if $\mathcal{A}_{p,i}$ contains a finite path from an initial state to an
accepting state. Let
${R}_{p}^i = q_{p,0}^i q_{p,1}^i \cdots q_{p,P}^i$
be such a finite path, where
$q_{p,0}^i \in Q_{p,i,0}$,
$q_{p,P}^i \in \mathcal{F}_{p,i}$, and
$(q_{p,\ell}^i,q_{p,\ell+1}^i) \in \delta_{p,i}$ for all
$\ell=0,1,\cdots,P-1$. Its cost is defined by
\begin{equation}
\label{ch03:eq:init-cost}
\texttt{Cost}(R_p^i,\mathcal{A}_{p,i}) \triangleq
\sum_{\ell=0}^{P-1} W_{p,i}(q_{p,\ell}^i,q_{p,\ell+1}^i),
\end{equation}
which is the total weight accumulated along $R_p^i$. The notation
$R_p^i[\ell]$ refers to the $\ell$th element of the path, and
$R_p^i[\ell\!:\!k]$ denotes the segment from the $\ell$th element to the $k$th
element.

\begin{problem}
\label{ch03:prob-accept}
Find a finite path $R_p^i$ of $\mathcal{A}_{p,i}$ that starts from an initial
state, ends at an accepting state, and minimizes the cost in
\eqref{ch03:eq:init-cost}.
\end{problem}

Let $R_{p,\textup{init}}^i$ solve Problem~\ref{ch03:prob-accept}.
Algorithm~\ref{ch02:alg:optimal} computes it by shortest-path search on the
product automaton \cite{guo2013motion,lavalle2006planning,korf1987planning},
with worst-case complexity
$\mathcal{O}(|\delta_{p,i}| \cdot \log |Q_{p,i}| \cdot |Q_{p,i,0}|)$.
Projection gives
$\tau_{\mathcal{R},\textup{init}}^i \triangleq
R_{p,\textup{init}}^i|_{\Pi_{\mathcal{R}}^i}$, which fulfills $\varphi_i$.
These plans are synthesized locally rather than by a central coordinator
\cite{chen2012formal,guo2012infeasible}. If
$\tau_{\mathcal{R},\textup{init}}^i$ contains collaborative actions, execution
requires online assistance from other robots.

(II) \emph{Plan in a finite horizon and request generation}.
Let $\pi_{\mathcal{R},t}^i \in \Pi_{\mathcal{R}}^i$ denote the current state of
robot $i$ at time $t$, and suppose that
$\pi_{\mathcal{R},t}^i=\tau_{\mathcal{R},\textup{init}}^i[\ell]$ for some
index $\ell$. Each robot is assigned a finite planning horizon
$0 < H_i < \infty$, which specifies how far ahead the robot examines its plan.
The segment of the initial plan expected to be reached within this horizon is
denoted by $\tau_{\mathcal{R},H}^i$ and is defined by
$\tau_{\mathcal{R},H}^i \triangleq
\tau_{\mathcal{R},\textup{init}}^i[\ell\!:\!k]$,
where $k \ge \ell$ is the smallest index satisfying
$
\sum_{s=\ell}^{k-1}
T_{\mathcal{R}}^i\!\left(
\tau_{\mathcal{R},\textup{init}}^i[s],
\tau_{\mathcal{R},\textup{init}}^i[s+1]
\right)
\ge H_i$.
If no such index exists, then $k$ is set to the final index of the remaining
plan. This horizon prevents coordination attempts for collaborative actions
that lie too far in the future.

The robot then checks whether $\tau_{\mathcal{R},H}^i$ contains a collaborative
action. Let $\langle \pi_v,\sigma_m \rangle$ be the first state in
$\tau_{\mathcal{R},H}^i$ such that $\sigma_m \in \Sigma_c^i$. In that case,
robot $i$ broadcasts a request to its neighbors. The request message is given by:
\begin{equation}
\label{ch03:request}
\textbf{Request}^i \triangleq
\{(\sigma_d,\pi_v,T_m) \mid
  \sigma_d \in \texttt{Depd}_i(\sigma_m)\},
\end{equation}
where $\pi_v$ is the region at which the collaborative action will be executed,
and $T_m \ge 0$ is the planned time from now until the execution of
$\sigma_m$. If $\langle \pi_v,\sigma_m \rangle$ is the $k$th state of
$\tau_{\mathcal{R},H}^i$, then
$T_m \triangleq
\sum_{s=\ell}^{k-1}
T_{\mathcal{R}}^i\!\left(
\tau_{\mathcal{R},H}^i[s],
\tau_{\mathcal{R},H}^i[s+1]
\right)$,
which is the elapsed time from the current state to the requested
collaboration. Thus $(\sigma_d,\pi_v,T_m)$ requests assisting action
$\sigma_d$ at region $\pi_v$ after time $T_m$, and the same message
$\textbf{Request}_j^i$ is sent to each neighbor. Only the first collaborative
action in the horizon is considered, because its outcome may alter later plan
segments.

%%%%%%%%%%%%%%%%%%%%%%%%%%%%%%%%%%%%%%%%%%%%%%%%%%%
\begin{algorithm}[t]
\caption{Plan in horizon and request, \texttt{Request}()}
\label{ch03:alg:request}
\DontPrintSemicolon

\KwIn{$\tau_{\mathcal{R},\textup{init}}^i$, $\pi_{\mathcal{R},t}^i$, $H_i$}
\KwOut{$\tau_{\mathcal{R},H}^i$, $\textbf{Request}^i$}

Find $\ell$ such that
$\tau_{\mathcal{R},\textup{init}}^i[\ell]=\pi_{\mathcal{R},t}^i$\;
$k=\ell$, $T_m=0$, $\textbf{Request}^i=\emptyset$\;

\While{$(T_m < H_i)$ and $(k+1 < |\tau_{\mathcal{R},\textup{init}}^i|)$}{
  $k=k+1$\;
  $T_m=T_m+
  T_{\mathcal{R}}^i\!\left(
  \tau_{\mathcal{R},\textup{init}}^i[k-1],
  \tau_{\mathcal{R},\textup{init}}^i[k]\right)$\;
  $\langle \pi_v,\sigma_m \rangle=
  \tau_{\mathcal{R},\textup{init}}^i[k]$\;

  \If{$\textup{\textbf{Request}}^i=\emptyset$ and $\sigma_m \in \Sigma_c^i$}{
    \ForAll{$\sigma_d \in \textup{\texttt{Depd}}_i(\sigma_m)$}{
      add $(\sigma_d,\pi_v,T_m)$ to $\textbf{Request}^i$\;
    }
  }
}

$\tau_{\mathcal{R},H}^i=
\tau_{\mathcal{R},\textup{init}}^i[\ell\!:\!k]$\;

\Return{$\tau_{\mathcal{R},H}^i$, $\textup{\textbf{Request}}^i$}\;
\end{algorithm}
%%%%%%%%%%%%%%%%%%%%%%%%%%%%%%%%%%%%%%%%%%%%%%%%%%%

(III) \emph{Request evaluation and reply}.
After receiving a request, a neighbor
$j \in \mathcal{N}_i$ evaluates it in terms of feasibility and timing. The
reply message from robot $j$ to robot $i$ is given by:
\begin{equation}
\label{ch03:eq:reply}
\textbf{Reply}_i^j \triangleq
\{(\sigma_d,b_d^j,t_d^j) \mid
  (\sigma_d,\pi_v,T_m) \in \textbf{Request}_j^i\},
\end{equation}
where $b_d^j \in \{\top,\bot\}$ indicates feasibility and $t_d^j \ge 0$ is the
available assistance time. Let $\overline{T}_j$ be the finishing time of the
collaboration currently assigned to robot $j$. If $\overline{T}_j>0$, robot
$j$ rejects all requests; otherwise it evaluates them. With current accepting
run $R_{p,t^-}^j$, current product state $q_{p,t}^j$, accepting state
$q_{p,f}^j$, and remaining segment
$\widehat{R}_{p,-}^j$, the objective is to find an alternative segment
$\widehat{R}_{p,+}^j$ that reaches $\langle \pi_v,\sigma_d \rangle$ near time
$T_m$ with limited additional cost.
Let
$S_d \triangleq
\{q_p \in Q_{p,j} \mid
  q_p|_{\Pi_{\mathcal{R}}^j}=\langle \pi_v,\sigma_d \rangle\}$,
which is the set of product states associated with the requested assisting
action. A candidate revised run has the structure below:
\begin{equation}
\label{ch03:segment}
\widehat{R}_{p,+}^j \triangleq
q_{p,t}^j \cdots q_{p,r}^j \cdots q_{p,f}^j,
\end{equation}
where $q_{p,r}^j \in S_d$. The balanced cost of
$\widehat{R}_{p,+}^j$ is defined by:
\begin{equation}
\label{ch03:balanced}
\begin{aligned}
\texttt{BalCost}(\widehat{R}_{p,+}^j,T_m,\mathcal{A}_{p,j})
\triangleq\;&
\left|
\sum_{s=1}^{r-t} W_{p,j}\big(\widehat{R}_{p,+}^j[s],\widehat{R}_{p,+}^j[s+1]\big)
- T_m
\right| \\
& + \alpha_j \big(
\texttt{Cost}(\widehat{R}_{p,+}^j,\mathcal{A}_{p,j})
-\texttt{Cost}(\widehat{R}_{p,-}^j,\mathcal{A}_{p,j})
\big),
\end{aligned}
\end{equation}
where the first term measures the timing mismatch relative to the request and
the second term measures the additional cost of the revised run. The parameter
$\alpha_j>0$ is a design weight.
The corresponding optimization problem is as follows.

%%%%%%%%%%%%%%%%%%%%%%%%%%%%%%%%%%%%%%%%%%%%%%%%%%%
\begin{algorithm}[t]
\caption{Reply to request by robot $j$, \texttt{Reply}()}
\label{ch03:alg:reply}
\DontPrintSemicolon

\KwIn{$\textbf{Request}_j^i$, $\widehat{R}_{p,-}^j$, $\mathcal{A}_{p,j}$,
$\overline{T}_j$}
\KwOut{$\textbf{Reply}_i^j$, $\widehat{P}$}

$\textbf{Reply}_i^j=\emptyset$, $\widehat{P}=\emptyset$\;

\If{$\overline{T}_j>0$}{
  \ForAll{$(\sigma_d,\pi_v,T_m) \in \textbf{Request}_j^i$}{
    add $(\sigma_d,\bot,0)$ to $\textbf{Reply}_i^j$\;
  }
  \Return{$\textbf{Reply}_i^j$, $\widehat{P}$}\;
}

\ForAll{$(\sigma_d,\pi_v,T_m) \in \textup{\textbf{Request}}_j^i$}{
  $(\widehat{R}_{p,+}^j,b_d^j,t_d^j)=
  \texttt{EvalReq}(
  \widehat{R}_{p,-}^j,(\pi_v,\sigma_d,T_m),\mathcal{A}_{p,j})$\;

  \If{$b_d^j=\top$}{
    $\widehat{P}(\sigma_d)=\widehat{R}_{p,+}^j$\;
    add $(\sigma_d,\top,t_d^j)$ to $\textbf{Reply}_i^j$\;
  }
  \Else{
    add $(\sigma_d,\bot,0)$ to $\textbf{Reply}_i^j$\;
  }
}

\Return{$\textup{\textbf{Reply}}_i^j$, $\widehat{P}$}\;
\end{algorithm}
%%%%%%%%%%%%%%%%%%%%%%%%%%%%%%%%%%%%%%%%%%%%%%%%%%%


\begin{problem}
Given $\widehat{R}_{p,-}^j$, $S_d$, and $\mathcal{A}_{p,j}$, find a path
segment $\widehat{R}_{p,+}^j$ with the form \eqref{ch03:segment} that
minimizes \eqref{ch03:balanced}.
\end{problem}



Algorithm~\ref{ch03:adapt} solves this problem by bidirectional Dijkstra
search. The routine \texttt{DijksTA}($\cdot$) computes shortest paths from a
source to a target set while avoiding specified states. Let
\begin{equation}
\label{ch03:eq:Sc}
S_c \triangleq
\{q_p \in Q_{p,j} \mid
  q_p|_{\Pi_{\mathcal{R}}^j}
  = \langle \pi_{\mathcal{R}}^j,\sigma_n \rangle,\,
  \sigma_n \in \Sigma_c^j \cup \Sigma_h^j\},
\end{equation}
which excludes other collaborative and assisting actions from the forward
search. The forward and reversed calls to \texttt{DijksTA} identify candidate
paths through $S_d$, and the state minimizing \eqref{ch03:balanced} determines
$\widehat{R}_{p,+}^j$. If no such state exists, robot $j$ cannot assist. The
complexity of \texttt{DijksTA} is
$\mathcal{O}(|\delta_{p,j}| \cdot \log |Q_{p,j}|)$, and reversing the automaton
is linear in $|\delta_{p,j}|$
\cite{baier2008principles,lavalle2006planning,korf1987planning}. The candidate
run is stored in $\widehat{P}$ until robot $i$ sends confirmation; simultaneous
requests are processed sequentially.

\begin{lemma}
\label{ch03:td}
If Algorithm~\ref{ch03:adapt} returns $b_d^j=\top$, then action $\sigma_d$ can
be done by robot $j$ at time $t_d^j$ by following the run
$\widehat{R}_{p,+}^j$.
\end{lemma}

\begin{proof}
The first segment of $\widehat{R}_{p,+}^j$, from $q_{p,t}^j$ to the selected
state $q_{p,r}^{j,\star}$, is produced by \texttt{DijksTA} while avoiding all
states in $S_c$. Hence this segment contains no collaborative or assisting
actions other than the requested action $\sigma_d$. It therefore consists only
of motions and local actions that robot $j$ can execute independently. The time
required to reach $q_{p,r}^{j,\star}$ is exactly $t_d^j$.
\end{proof}


(IV) \emph{Confirmation and assignment}.
Based on the replies from the neighbors, robot $i$ sends confirmation messages
of the form:
\begin{equation}
\label{ch03:eq:confirm}
\textbf{Confirm}_j^i \triangleq
\{(\sigma_d,c_d^j,f_m) \mid
  \sigma_d \in \texttt{Depd}_i(\sigma_m)\},
\end{equation}
where $c_d^j \in \{\top,\bot\}$ indicates confirmation and $f_m$ is the
scheduled finishing time. Each assisting action must be assigned to exactly one
robot, each neighbor may support at most one assisting action for this
collaboration, and the finishing time is minimized. Let
$|\mathcal{N}_i|=N_1$, $|\texttt{Depd}_i(\sigma_m)|=N_2$,
$\mathcal{N}_i=\{1,\cdots,N_1\}$, and
$\texttt{Depd}_i(\sigma_m)=\{\sigma_1,\cdots,\sigma_{N_2}\}$. The assignment is
\begin{equation}
\label{ch03:integer}
\begin{aligned}
\textbf{minimize} \quad & f_m, \\
\textbf{subject to} \quad
& f_m \ge T_m, \\
& f_m \ge c_d^j t_d^j,
  \quad \forall j \in \{1,\cdots,N_1\},\;
  \forall d \in \{1,\cdots,N_2\}, \\
& \sum_{d=1}^{N_2} b_d^j c_d^j \le 1,
  \quad \forall j \in \{1,\cdots,N_1\}, \\
& \sum_{j=1}^{N_1} b_d^j c_d^j = 1,
  \quad \forall d \in \{1,\cdots,N_2\}, \\
& c_d^j \in \{0,1\},
  \quad \forall j \in \{1,\cdots,N_1\},\;
  \forall d \in \{1,\cdots,N_2\},
\end{aligned}
\end{equation}
where $(\sigma_d,b_d^j,t_d^j) \in \textbf{Reply}_i^j$ from
\eqref{ch03:eq:reply}. Any standard integer-programming solver can be used once
\eqref{ch03:integer} has been formed.
If \eqref{ch03:integer} is feasible, robot $i$ confirms the selected helper for
each assisting action and rejects the others. If it is infeasible, all helpers
are rejected and $\sigma_m$ must be delayed. The optimization is local to robot
$i$ and contains
$|\mathcal{N}_i| \cdot |\texttt{Depd}_i(\sigma_m)|$ Boolean variables.

%%%%%%%%%%%%%%%%%%%%%%%%%%%%%%%%%%%%%%%%%%%%%%%%%%%
\begin{algorithm}[t]
\caption{Evaluate the request, \texttt{EvalReq}($\cdot$)}
\label{ch03:adapt}
\DontPrintSemicolon

\KwIn{$\widehat{R}_{p,-}^j$, $(\pi_v,\sigma_d,T_m)$, $\mathcal{A}_{p,j}$}
\KwOut{$\widehat{R}_{p,+}^j$, $b_d^j$, $t_d^j$}

$q_{p,t}^j=\widehat{R}_{p,-}^j[1]$,
$q_{p,f}^j=\widehat{R}_{p,-}^j[\text{-1}]$\;

Compute $S_d$ and $S_c$\;

$\bar{c}=\texttt{Cost}(\widehat{R}_{p,-}^j,\mathcal{A}_{p,j})$\;

$(P_1,C_1)=\texttt{DijksTA}
(\mathcal{A}_{p,j},q_{p,t}^j,S_d,S_c)$\;

$(P_2,C_2)=\texttt{DijksTA}
(\texttt{Reverse}(\mathcal{A}_{p,j}),q_{p,f}^j,S_d,\emptyset)$\;

\ForAll{$q_{p,r}^j \in S_d$}{
  \If{$P_1(q_{p,r}^j)$ and $P_2(q_{p,r}^j)$ exist}{
    $C_3(q_{p,r}^j)=
    |C_1(q_{p,r}^j)-T_m|
    +\alpha_j(C_1(q_{p,r}^j)+C_2(q_{p,r}^j)-\bar{c})$\;
  }
}

Find $q_{p,r}^{j,\star} \in S_d$ that minimizes $C_3(q_{p,r}^j)$\;

\If{$q_{p,r}^{j,\star} \neq \emptyset$}{
  $P=P_1(q_{p,r}^{j,\star})+
  \texttt{Reverse}(P_2(q_{p,r}^{j,\star}))$\;
  \Return{$P$, $\top$, $C_1(q_{p,r}^{j,\star})$}\;
}

\Return{$\emptyset$, $\bot$, $0$}\;
\end{algorithm}
%%%%%%%%%%%%%%%%%%%%%%%%%%%%%%%%%%%%%%%%%%%%%%%%%%%

(V) \emph{Plan adaptation}.
After confirmation, robot $i$ either executes $\sigma_m$ as planned and sets
$\overline{T}_i=f_m$, or, if \eqref{ch03:integer} is infeasible, delays
$\sigma_m$ by Algorithm~\ref{ch03:alg:delay}. The delayed segment reaches
$\langle \pi_v,\sigma_m \rangle$ no earlier than $T_m+D_i$; waiting is enabled
by the idle action $\sigma_0$. A confirmed neighbor replaces its current
segment with the stored run in $\widehat{P}$ and sets $\overline{T}_j=f_m$;
otherwise it keeps its plan. After completing its own finite task, a robot sets
$\varphi_i \triangleq \top$ and remains available for assistance.


\begin{theorem}
\label{ch03:tm}
If \eqref{ch03:integer} is feasible, the collaborative action $\sigma_m$
is accomplished at time $f_m$.
\end{theorem}

\begin{proof}
For each assisting action
$\sigma_d \in \texttt{Depd}_i(\sigma_m)$, feasibility of
\eqref{ch03:integer} implies that exactly one robot
$j \in \mathcal{N}_i$ is assigned to $\sigma_d$. By
Lemma~\ref{ch03:td}, that robot can complete $\sigma_d$ at time $t_d^j$. By
the constraints in \eqref{ch03:integer}, $f_m$ is at least the scheduled time
of every confirmed assisting action and is no earlier than $T_m$. Hence all
assisting actions required by $\sigma_m$ are completed by time $f_m$, and
$\sigma_m$ is accomplished at time $f_m$.
\end{proof}


(VI) \emph{Failure detection and recovery}.
For collaborative tasks, robot $i$ monitors confirmed helpers until
$\sigma_m$ is completed. If a helper $h \in \mathcal{N}_i$ fails to acknowledge
within a bounded time, it is treated as failed and the affected assisting action
$\sigma_d$ is reassigned. Robot $i$ sends a new request and selects the
replacement by
\begin{equation}
\label{ch03:eq:reassign}
\hat{h} \triangleq
\arg\min_{j \in \mathcal{N}_i}
|f_m - t_d^j b_d^j|,
\end{equation}
where $f_m$ is the previously confirmed finishing time. Third, it sends the
confirmation $\{(\sigma_d,\top,f_m)\}$ to robot $\hat{h}$. Finally, robot
$\hat{h}$ adapts its plan according to the nominal procedure.


%%%%%%%%%%%%%%%%%%%%%%%%%%%%%%%%%%%%%%%%%%%%%%%%%%%
\begin{algorithm}[t!]
\caption{Delay collaboration, \texttt{DelayCol}()}
\label{ch03:alg:delay}
\DontPrintSemicolon

\KwIn{$\widehat{R}_{p,-}^i$, $(\pi_v,\sigma_m,T_m)$, $D_i$, $\mathcal{A}_{p,i}$}
\KwOut{$\widehat{R}_{p,+}^i$}

$q_{p,t}^i=\widehat{R}_{p,-}^i[1]$,
$q_{p,f}^i=\widehat{R}_{p,-}^i[\text{-1}]$\;

Compute
$S_d=
\{q_p \in Q_{p,i} \mid
  q_p|_{\Pi_{\mathcal{R}}^i}
  = \langle \pi_v,\sigma_m \rangle\}$ and $S_c$\;

$(P_1,C_1)=\texttt{DijksTA}
(\mathcal{A}_{p,i},q_{p,t}^i,S_d,S_c)$\;

$(P_2,C_2)=\texttt{DijksTA}
(\texttt{Reverse}(\mathcal{A}_{p,i}),q_{p,f}^i,S_d,\emptyset)$\;

\ForAll{$q_{p,d}^i \in S_d$}{
  \If{$C_1(q_{p,d}^i)>T_m+D_i$ and $P_2(q_{p,d}^i)$ exists}{
    $\widehat{R}_{p,+}^i=
    P_1(q_{p,d}^i)+
    \texttt{Reverse}(P_2(q_{p,d}^i))$\;
    \Return{$\widehat{R}_{p,+}^i$}\;
  }
}
\end{algorithm}
%%%%%%%%%%%%%%%%%%%%%%%%%%%%%%%%%%%%%%%%%%%%%%%%%%%


\subsubsection{Overall structure}
\label{ch03:structure}

The coordination scheme is intended for loosely coupled multi-robot
systems, in which collaborations are local and relatively sparse compared with
the total number of robot activities. This property is formalized by the
following assumption.
Assumption~\ref{ch03:assump:loose} below states that every collaborative action
needed for a local task can eventually be supported in finite time. During
execution, each robot either evaluates incoming requests and adapts after
confirmation, or examines its own plan, sends requests, solves
\eqref{ch03:integer}, and delays the collaboration by
Algorithm~\ref{ch03:alg:delay} when needed. The correctness statement follows.

\begin{assumption}
\label{ch03:assump:loose}
There exists a finite constant ${T}>0$ such that, for every robot
$i \in \mathcal{N}$ and every collaborative action $\sigma_m$ requested by
robot $i$ at time $t_m>0$, the integer program in~\eqref{ch03:integer}
becomes feasible no
later than time $t_m+{T}$.
\end{assumption}

\begin{theorem}
\label{ch03:all}
Under Assumption~\ref{ch03:assump:loose}, the proposed coordination scheme
solves Problem~\ref{ch03:big-problem}. That is, all local tasks $\varphi_i$ are
fulfilled in finite time for all $i \in \mathcal{N}$.
\end{theorem}

\begin{proof}
Consider an arbitrary robot $i \in \mathcal{N}$ and its initial plan
$\tau_{\mathcal{R},\textup{init}}^i$. All motion transitions and local actions
in this plan can be executed autonomously by robot $i$. It remains to show that
each collaborative action required by the plan is also completed in finite
time.
If the plan of robot $i$ remains unchanged, then every collaborative action
$\sigma_m$ appearing in the plan is associated with a feasible instance of
\eqref{ch03:integer}. By Theorem~\ref{ch03:tm}, each such action is completed
at a finite time $f_m$. Since the plan is finite, the entire plan is executed
in finite time, and the corresponding trace satisfies $\varphi_i$.
If the plan of robot $i$ is revised during execution, two cases may occur.
First, robot $i$ may be confirmed to assist a neighboring robot. In that case,
Algorithm~\ref{ch03:adapt} constructs a revised accepting run that still
reaches an accepting state, so the updated plan continues to satisfy
$\varphi_i$ in finite time. Second, robot $i$ may request assistance for a
collaborative action $\sigma_m$ and receive an infeasible outcome for
\eqref{ch03:integer}, which triggers Algorithm~\ref{ch03:alg:delay}. By
Assumption~\ref{ch03:assump:loose}, the delayed coordination problem becomes
feasible within at most ${T}$ additional time, so the postponed
collaborative action is still completed in finite time.
Therefore, every robot fulfills its local task in finite time.
\end{proof}
